% A representative ACM acmart-style conference paper. Written for ByeTex
% test purposes; uses the documented acmart patterns without requiring the
% acmart.cls file to be present locally. Modeled after the official
% sample-sigconf.tex in https://github.com/borisveytsman/acmart.
\documentclass[sigconf]{acmart}

\usepackage{booktabs}

\setcopyright{acmlicensed}
\copyrightyear{2026}
\acmYear{2026}
\acmConference[Conference '26]{ACM Conference}{June 03--05, 2026}{Woodstock, NY}

\begin{document}

\title{Hello, World: A Sample ACM SIGCONF Paper}

\author{Alice Anderson}
\authornote{Both authors contributed equally.}
\affiliation{%
  \institution{Example University}
  \city{Springfield}
  \country{USA}
}
\email{alice@example.edu}

\author{Bob Brown}
\affiliation{%
  \institution{Example University}
  \country{USA}
}
\email{bob@example.edu}

\renewcommand{\shortauthors}{Anderson and Brown}

\begin{abstract}
This paper demonstrates a small ACM-formatted submission. We show common
constructs that ByeTex must handle: sectioning, math, tables, and citations
in the natbib-style used by \LaTeX{} authors.
\end{abstract}

\keywords{datasets, neural networks, gaze detection, text tagging}

\maketitle

\section{Introduction}

Following standard \LaTeX{} practice, sections begin with \verb|\section{...}|.
Tools like \BibTeX{} format the bibliography automatically.

Citations come in two flavors~\cite{bohr}: numerical or author-year.

\section{Methodology}

The fundamental relation is given by Einstein:
\begin{equation}
E = mc^2.
\label{eq:einstein}
\end{equation}

We also use the canonical Pythagorean identity, $\sin^2 x + \cos^2 x = 1$, and
the change-of-variables form $\int_0^1 x^2 \, dx = \frac{1}{3}$.

\subsection{Setup}

Table~\ref{tab:results} summarizes our experimental results.

\begin{table}[h]
\caption{Sample experimental results.}
\label{tab:results}
\begin{tabular}{lcr}
\toprule
Method & Accuracy & Time (s) \\
\midrule
Baseline & 0.81 & 12.3 \\
Ours     & 0.93 & 14.7 \\
\bottomrule
\end{tabular}
\end{table}

\section{Results}

As shown in Equation~\eqref{eq:einstein}, energy and mass relate by the speed
of light squared. Our improvements over the baseline are 12 percentage points.

\section{Conclusion}

This template demonstrates the bones of an ACM SIGCONF paper. Refer to
\cite{einstein,dirac} for additional context.

\bibliographystyle{ACM-Reference-Format}
\bibliography{sample-base}

\end{document}
